\section{Introduction}\label{sec:intro}
The three-dimensional incompressible Navier--Stokes equations with
viscosity $\nu>0$ and external force $f$ are
\begin{equation}\label{eq:NS}
 \partial_tu+(u\cdot\nabla)u-\nu\Delta u+\nabla p=f,
 \qquad \nabla\cdot u=0,\qquad u(\cdot,0)=a.
\end{equation}
Here $u$ is the velocity, $p$ is the pressure, and $a$ is the prescribed
initial velocity. We work on $D=\R^3$ or on the unit torus
$D=\TT=\R^3/\mathbb Z^3$, with volume one and periodic velocity and
pressure in the latter case. Our starting point is the following theorem, stated with
the notation of OpenAI~\cite[Theorem~1.1]{openai}.

\Needspace{12\baselineskip}
\begin{theorem}[OpenAI]\label{thm:packet}
For every $\nu>0$ there exist a force
$f\in C_c^\infty(\R^3\times(0,\infty);\R^3)$, a compact set
$K\subset\R^3$, and smooth velocity and pressure fields $u,p$ on
$\R^3\times[0,1)$ satisfying
\begin{equation}\label{eq:packet}
 \partial_tu+(u\cdot\nabla)u-\nu\Delta u+\nabla p=f,
 \qquad \nabla\cdot u=0,\qquad u(\cdot,0)=0,
\end{equation}
such that $\supp u(\cdot,t)\cup\supp p(\cdot,t)\subset K$ for every
$0\le t<1$,
\begin{equation}\label{eq:packetblowup}
 \sup_{0\le t<1}\norm{u(t)}_{L^2(\R^3)}<\infty,
 \qquad \limsup_{t\uparrow1}\norm{u(t)}_{L^\infty(\R^3)}=\infty.
\end{equation}
Consequently, there is no smooth solution $(u,P)$ on
$\R^3\times[0,\infty)$ with the same force and initial datum whose
kinetic energy is uniformly bounded
$\sup_{t\ge0}\frac12\int_{\R^3}|u(x,t)|^2\dd x<\infty$.
\end{theorem}

Fix a viscosity $\nu>0$, an initial velocity $a$, and a time $T>0$.
Our aim is to determine in which norms smooth forces producing breakdown
of the classical Navier--Stokes solution by time $T$ are dense among smooth
forces. The force is allowed to vary, while the initial velocity is kept
fixed. This formulation asks how closely a regular flow can be approximated
by a flow that becomes singular, when closeness is measured in a prescribed
force norm.

Our construction follows a local insertion strategy closely related to
the work of Enciso, Pe\~nafiel-Tom\'as and Peralta-Salas~\cite{eppt2025}
on the incompressible Euler equations. Their Theorem~1.5 constructs weak
solutions that agree with a prescribed smooth flow outside a chosen region
and have singular behavior inside it. Lemma~2.11 and Section~11 combine
local modification of the background, preservation of incompressibility,
and insertion of localized singular solutions. In particular,
Section~11.3 first makes the background spatially constant in a small
region before inserting the singular construction. Our argument uses
the same general strategy, with the compact classical forced
Navier--Stokes solution of Theorem~\ref{thm:packet} as its input.
Here the smooth gluing defect is absorbed into a modified external force;
in~\cite{eppt2025}, convex integration removes the Reynolds stress to
produce unforced weak Euler solutions. We study the resulting quantitative
force-density consequences, preserving the initial velocity and earlier
history and identifying the Sobolev thresholds stated below.

In Theorem~\ref{thm:packet}, $C_c^\infty$ denotes smooth functions with compact support in the
indicated open set, so the force is smooth through the blowup time and
vanishes near $t=0$. For the construction below, fix one solution from
Theorem~\ref{thm:packet} and denote its velocity, pressure, and force by
$(U,P,F)$. This distinguishes it from the reference solution to be
modified. The corresponding compact solution and its support and
zero-data properties are also recorded in the source repository~\cite{lean}.

We now define the spaces in which the force approximation is measured.
For background on Fourier transforms, distributions, and Sobolev spaces,
see Taylor~\cite[Chapter~3, Sections~3--4, and Chapter~4, Sections~1 and~3]{taylor2011};
the definitions and normalizations used here are specified explicitly
below.\footnote{On $\R^3$, our Fourier normalization and inhomogeneous
Sobolev norm agree with Taylor's. His Bessel multiplier $\Lambda^s$
corresponds to our $J^s=(I-\Delta)^{s/2}$, whereas we use
$\Lambda^s=(-\Delta)^{s/2}$. His torus formulas use $2\pi$-periodic
angular coordinates; we use the unit torus. The weights
$(1+|k|^2)^s$ and $(1+4\pi^2|k|^2)^s$ give equivalent, not identical,
periodic norms after the coordinate change.}
For $s\in\R$ and $D\in\{\R^3,\TT\}$, $H^s(D)$ denotes the
inhomogeneous $L^2$-based Sobolev space of order $s$. Its norm is
\begin{align*}
 \norm{z}_{H^s(\TT)}^2
 &=\sum_{k\in\mathbb Z^3}(1+4\pi^2|k|^2)^s|\widehat z(k)|^2,\\
 \norm{z}_{H^s(\R^3)}^2
 &=\int_{\R^3}(1+|\xi|^2)^s|\widehat z(\xi)|^2\dd\xi,
\end{align*}
where
$\widehat z(k)=\int_{\TT}z(x)e^{-2\pi i k\cdot x}\dd x$ on the torus
and
$\widehat z(\xi)=(2\pi)^{-3/2}\int_{\R^3}e^{-ix\cdot\xi}z(x)\dd x$
on the whole space. More precisely,
\begin{align*}
 H^s(\R^3)&=\{z\in\mathcal S'(\R^3):
        (1+|\xi|^2)^{s/2}\widehat z\in L^2(\R^3)\},\\
 H^s(\TT)&=\{z\in\mathcal D'(\TT):
        \norm{z}_{H^s(\TT)}<\infty\}.
\end{align*}
Here $\mathcal S'(\R^3)$ is the space of tempered distributions and
$\mathcal D'(\TT)$ the space of periodic distributions; Fourier transforms
and coefficients are understood distributionally. Thus $H^0(D)=L^2(D)$,
and negative orders allow distributions that need not be $L^2$ functions.
For vectors and tensors we sum the squared component norms.
The notation $\norm{z}_p$ means the spatial $L^p(D)$
norm. The homogeneous norm $\dot H^s$ replaces the weights above by
$(2\pi|k|)^{2s}$ or $|\xi|^{2s}$, respectively. On the torus we omit
$k=0$ and require zero mean; on $\R^3$ the homogeneous norm is used
when finite, with the relevant distributional realizations specified
in Section~\ref{sec:preliminaries} and Appendix~\ref{sec:critical-appendix}.
For background on homogeneous Sobolev norms, see
Tao~\cite[Appendix~A]{taodispersive}.\footnote{Tao uses the unnormalized
Fourier transform and includes a compensating factor $(2\pi)^{-3/2}$
in the Fourier formula for the Sobolev norm on $\R^3$; the resulting
norm agrees with ours wherever finite. This reference does not replace
our explicit conventions for low frequencies and homogeneous-space
realizations.}

For a time interval $I$ and a spatial normed space $X$, $L^q(I;X)$ is the
space of strongly measurable maps with finite norm; for the Bochner-space
framework, see Hunter~\cite[Appendix~6.A, Definitions~6.14 and~6.27]{hunterpde}.\footnote{Hunter
states these definitions for Banach-valued functions on a finite interval
$(0,T)$. Here the same measurability condition and norm formulas are used
on the indicated interval $I$, including $(0,\infty)$; functions equal
almost everywhere are identified.}
\begin{equation}\label{eq:time-norms}
 \norm{z}_{L^q(I;X)}=\left(\int_I\norm{z(t)}_X^q\dd t\right)^{1/q},
 \qquad
 \norm{z}_{L^\infty(I;X)}=\operatorname*{ess\,sup}_{t\in I}\norm{z(t)}_X,
\end{equation}
where the first formula applies for $1\le q<\infty$. In particular,
\[
 L^q_tH^s_x=L^q(0,\infty;H^s(D)),\qquad
 L^q_t\dot H^s_x=L^q(0,\infty;\dot H^s(D)),\qquad
 L^q_tL^p_x=L^q(0,\infty;L^p(D))
\]
for forces: the time exponent controls integrability, and the Sobolev
order measures only spatial regularity. We give each smooth force class
the relative topology induced by the stated norm. This differs from the
test-function topology, which controls all derivatives on a common compact
support. Unless an interval is displayed, force norms use $(0,\infty)$
and velocity norms before blowup use $(0,T)$. For a fixed finite time
$T>0$, the velocity approximation is measured by
\begin{equation}\label{eq:Enorm}
 \norm{z}_{E_T}=\norm{z}_{L^\infty(0,T;L^2(D))}
                +\norm{\nabla z}_{L^2(0,T;L^2(D))}.
\end{equation}
Since $\norm{z}_{L^2_tL^2_x}\le T^{1/2}\norm{z}_{L^\infty_tL^2_x}$,
this is equivalent to the sum norm on
$L^\infty_tL^2_x\cap L^2_tH^1_x$ on a fixed finite interval; it imposes
no endpoint value at $T$.

The construction combines compact support, parabolic scaling, and smooth
cutoffs. Scaling places the singular solution in an arbitrarily small
spatial ball and a short interval ending at $T$. To insert it into a
regular reference solution $v$, one must also remove $v$ near that ball:
direct addition would create nonlinear interaction terms involving the
singular velocity. The cutoff construction is the smooth form of the
Urysohn lemma: a function is chosen to equal one near a compact set and
zero outside a slightly larger open set. Incompressibility requires a
further step. Locally write $v=\nabla\times A$ and set
\[
 w_\eps=-\nabla\times(\eta_\eps\theta_\eps A),
 \qquad u_\eps=v+w_\eps+U_\eps,
\]
where $U_\eps$ is the rescaled singular solution. The cutoffs are chosen
so that $v+w_\eps$ vanishes on a neighborhood of its active support.
Both nonlinear interaction terms then vanish identically. The new force
is the sum of the original force, the rescaled smooth force, and a smooth
correction determined entirely by $v$ and the cutoffs.

This construction explains the positive density results. A force rescales
with amplitude $\eps^{-3}$, spatial volume $\eps^3$, and time scale
$\eps^2$. Its homogeneous $L^q_t\dot H^s_x$ scaling factor, whenever that
norm is finite, is
\[
 \eps^{\,2/q-3/2-s}.
\]
The background correction has one additional power of $\eps$. On the
whole space, negative orders require a separate estimate near frequency
zero. On the torus, a localization estimate compares a compactly supported
function with its periodization. These arguments give the following
thresholds for breakdown from rest:
\begin{center}
\begin{tabular}{lll}
\toprule
Domain & Force topology & Density holds exactly when\\
\midrule
$\TT$ & $L^1(0,\infty;H^s)$ & $s<1/2$\\
$\R^3$ & $L^1(0,\infty;H^s)$ & $s<1/2$\\
$\R^3$ & $L^2(0,\infty;H^s)$ & $s<-1/2$\\
\bottomrule
\end{tabular}
\end{center}
Theorems~\ref{thm:main} and~\ref{thm:Rmain} give the precise statements.
Below each threshold, density holds for every fixed admissible smooth
initial velocity. The converse statements are proved for initial velocity
zero, by constructing a nonempty open set of regular forces. In particular,
they require estimates beyond the scaling calculation.

The approximation is also strong at the level of the velocity. Its
$L^\infty_tL^2_x$ difference and the $L^2_{t,x}$ norm of its gradient both
tend to zero. The solution agrees with the reference until shortly before
$T$, and its speed becomes unbounded exactly at $T$ whenever the reference
is regular through that time. For density around a general force, the
conclusion is breakdown \emph{by} $T$: a force whose solution already
breaks down earlier is left unchanged.

A related force-density result was proved by Hofmanov\'a, Zhu and Zhu
\cite[Theorem~1.4, Corollary~4.6, and Theorem~5.1]{hzz2024}. In the
deterministic setting, forces yielding nonunique Leray--Hopf solutions
from rest are dense in $L^1(0,1;L^2(\R^3))$, and each divergence-free
$L^2$ initial velocity admits a force yielding nonuniqueness. Here the
property under study is classical breakdown with a globally smooth force.
The compact blowup theorem supplies that property, while the local
construction above places it near an arbitrary regular reference.
Fujita and Kato~\cite{fujita} and Tao~\cite{tao2018} provide background
for the critical velocity framework. We use the classical forced local
theory of Tao~\cite[Theorems~5.1 and~5.4]{tao2013} and the classical
Sobolev embeddings cited in the appendices. The precise continuation
criterion, normalization reductions, and critical force estimates needed
here are justified below.

Section~\ref{sec:preliminaries} specifies the smooth data classes,
pressure conventions, and local theory, and derives an energy estimate
for the compact singular solution. Section~\ref{sec:torus} proves the periodic
result and its further consequences. Section~\ref{sec:whole} treats the
whole space, including rapidly decaying forces, completed force spaces,
and cell averages. Section~\ref{sec:discussion} discusses the scope of
the conclusions. Appendices~\ref{sec:local-appendix}
and~\ref{sec:critical-appendix} state the classical local theory and
Sobolev embeddings, with the continuation argument and the
normalization and homogeneous-space details needed on both domains.
